% concatenated from Sayan__1_.tex
\PassOptionsToPackage{dvipsnames}{xcolor}
\documentclass{aic}

\usepackage{amsmath,amsthm,amssymb,mathrsfs,amstext,enumitem, hyperref}
\usepackage[colorlinks=true]{hyperref}
\newtheorem{thm}{Theorem}
\newtheorem{lem}{Lemma}
\newtheorem{defn}{Definition}
\newcommand{\N}{\mathbb{N}}
\newcommand{\Z}{\mathbb{Z}}


\aicAUTHORdetails{%
  title = {Monochromatic Translated Product and Answering Sahasrabudhe's Conjecture}, %% please capitalize all significant words
  author = {Sayan Goswami},
    %% Please use the format for commas as follows:
    %% "A", or "A and B", or "A, B, and C", or "A, B, C, and D", etc.
  plaintextauthor = {Sayan Goswami},
    %% An author list in plain text: Use the format
    %% "A", or "A, B", or "A, B, C", etc.
    %% NOTE: No LaTeX code in author names.
    %% NOTE: No "and" at the end--simply comma separated,
    % 
 %% The remaining lines in this section are optional:
    %
    %% IF YOUR TITLE CONTAINS MATH OR LATEX such as accented characters: 
    %% Add a "plain text title";  otherwise comment out the next line:
  plaintexttitle = {Monochromatic Translated Product and Answering Sahasrabudhe's Conjecture}, %%  title without math or LaTeX
    %
    %% ONLY IF YOUR TITLE IS TOO LONG to fit in the page headers, please 
    %% add an abbreviated version of the title, otherwise comment it out:
  runningtitle = {Monochromatic Translated Product}, 
    %
    %% ONLY IF YOUR AUTHOR LIST IS TOO LONG to fit in the page headers, 
    %% add an abbreviated version, otherwise comment it out:
  runningauthor = {Sayan Goswami},
    %% you can replace first names and/or middle names with initials.
    %
    %% ONLY IF YOUR AUTHOR LIST IS TOO LONG to fit the copyright entry
    %% on the bottom of the front page,
    %% add an abbreviated version, otherwise comment it out:
  copyrightauthor = {Sayan Goswami},
    %% Note that the copyrightauthor  field will seldom be necessary;
    %% for instance, in this example with four authors, it would be 
    %% all right to comment it out and have all authors' full names 
    %% appear on the Copyright line
   %
   %% Include keywords of your choice: comma separated, lower case;
   %% comment out the "keywords" line if you don't wish to provide them
  keywords = {Schur Theorem, $IP$ sets, Hindman Conjecture, Moreira's sum product theorem.},
}   %%% END \aicAUTHORdetails

%%%%%%%%%%%%%%%%%%%%%%%%%%%%%%%%%%%%%%%%%%%%%%%%
%%% EDITOR: please fill in the following data:
\aicEDITORdetails{%
   year={2026},
   volume={},
   number={2},
   received={21 December 2024},   % received date: example: 7 January 2017
   revised={11 October 2025},    % Optional revised date (you may comment it out)
   published={10 April 2026},  % published date
   doi={10.19086/aic.2026.2},      % XXX = number of paper, e.g. aic006 for paper#6
%                              % or  aic0006 (length of string arbitrary)
}   %%% END \aicEDITORdetails







\begin{document}

\begin{frontmatter}[classification=text]
%% EDITOR: this will force the keywords to appear right after the Abstract.
%%   If the abstract is too long and would force the keywords off the
%%   front page, please comment out % [classification=text] above
%%   This way the keywords will be floated on the bottom of the first page
%%   even though the Abstract spills over to the next page.

%%% AUTHOR: Title goes here.  This line is optional.  You must use it
%%   if title has footnote attached or requires nontrivial typesetting,
%%   e.g., inclusion of linebreaks to force nice layout.
\title{Monochromatic Translated Product and Answering Sahasrabudhe's Conjecture} %% please capitalize all significant words

%%% AUTHOR:
%%% List all authors. If you wish, place grant acknowledgements in \thanks.
%%% In brackets include a short tag for each author.
\author[SG]{Sayan Goswami\thanks{Supported by  NBHM postdoctoral fellowship with reference no: 0204/27/(27)/2023/R \& D-II/11927.}}
%\author[johan]{Johan H{\aa}stad\thanks{Supported by...}}
%\author[laci]{L\'aszl\'o Lov\'asz\thanks{Supported by...}}
%\author[andy]{Andrew Chi-Chih Yao\thanks{Supported by...}}

%%% AUTHOR: Abstract goes here
\begin{abstract}
This article resolves two related problems in Ramsey theory on the integers. We show that for any finite coloring of the set of natural numbers, there exist numbers $a$ and $b$ for which the configuration $\{a, b, ab, a(b+1)\}$ is monochromatic. By redefining the variables $a=x$ and $ab=y,$ our configurations transforms into $\{x,y,x+y,\frac{y}{x}\}.$ This finding has two main consequences: first, it disproves a conjecture proposed by J. Sahasrabudhe; second, it establishes a quotient version of the long-standing Hindman's conjecture, which asks for a monochromatic set of the form $\{x,y,x+y,xy\}$. 
\end{abstract}
\end{frontmatter}

%%% AUTHOR: body of paper starts here
\section{Introduction}


Arithmetic Ramsey theory studies the existence of monochromatic configurations within any finite coloring of the integers or the natural numbers \( \mathbb{N} \). Here, a ``coloring'' refers to a finite disjoint partition of \( \mathbb{N} \), and a set is said to be \emph{monochromatic} if it is entirely contained in one part of this partition. A collection \( \mathcal{F} \) of subsets of \( \mathbb{N} \) is said to be \emph{partition regular}\footnote{Although this notion is often referred to as \emph{weakly partition regular} in the literature, we omit the term ``weakly'' throughout this paper.} if, for every finite coloring of \( \mathbb{N} \), there exists a monochromatic member \( F \in \mathcal{F} \).





In 1916, I.~Schur~\cite{schur} proved that the pattern \( \{a, b, a+b\} \) is partition regular. Applying the map \( n \mapsto 2^n \), one immediately obtains that \( \{a, b, a \cdot b\} \) is also partition regular. These two results are commonly referred to as the \emph{additive Schur theorem} and the \emph{multiplicative Schur theorem}, respectively.  



For more than a century, it remained unknown whether these two phenomena could be unified in a single statement. That is, whether \( \{a, b, a+b, a \cdot b\} \) is also partition regular (see~\cite{erdos}).  In 1979, R.~L.~Graham and N.~Hindman independently established such a result for two-colorings (see~\cite[Section~4]{trans}). Subsequently, N.~Hindman conjectured that the statement should hold for all finite colorings, a problem now known as the \emph{Hindman conjecture}. The strongest known partial result toward this conjecture is due to J.~Moreira~\cite{m}, who proved that there exist \( a, b \in \mathbb{N} \) such that the set \( \{a, ab, a+b\} \) is monochromatic; see also~\cite{al1} for a concise proof.  



In 2017, while investigating asymmetric variants of Schur’s theorem, J.~Sahasrabudhe~( see \cite[Question 31.]{saha}) conjectured that there exist finite colorings of the natural numbers that avoid monochromatic configurations of the form \( \{a, b, a(b+1)\} \). In this paper, we disprove this conjecture by showing that every finite coloring of \( \mathbb{N} \) necessarily contains such a configuration. In fact, we establish a stronger result: the larger set \( \{a, b, ab, a(b+1)\} \) can always be found monochromatic.  


On the other hand, by redefining the variables $a=x, ab=y$ our result shows that the pattern $\{x,y,x+y,\frac{y}{x}\}$ is monochromatic.
Our result can therefore be viewed as a quotient version of the Hindman conjecture. The technique of our proof builds upon the argument of Moreira~\cite{m}, who employed van der Waerden’s theorem, whereas we make iterative use of Hindman’s theorem.  



The following statement is our main result, which will be proved in the next section.

\begin{thm}\label{p}
For any finite coloring of \( \mathbb{N} \), there exist \( a, b \in \mathbb{N} \) such that the set \( \{a, b, ab, a(b+1)\} \) is monochromatic.
\end{thm}



%\section{Proof of Our Results}

\section{Proof of Theorem \ref{p}}

Before we proceed to the proof of our Theorem \ref{p}, we need to recall some technical definitions.

\begin{defn}\text{}
    \begin{enumerate}
        \item For any non-empty set $X,$ denote by $\mathcal{P}_f(X)$ the set of all non-empty finite subsets of $X.$
        \item For any sequence $\langle x_n\rangle_n$, denote by $$FS\left(\langle x_n\rangle_n\right)=\left\lbrace \sum_{n\in H}x_n:H\in \mathcal{P}_f(\N) \right\rbrace.$$

        \item A set $A\subseteq \N$ is called an IP set if there exists a sequence $\langle x_n\rangle_n$ such that $FS\left(\langle x_n\rangle_n\right)\subseteq A.$

         \item For any $R\in \N,$ a set $A\subseteq \N$ is called an $IP_R$ set if there exists a sequence $\langle x_n\rangle_{n=1}^R$ such that $FS\left(\langle x_n\rangle_{n=1}^R\right)\subseteq A.$

         \item For any $y\in \N,$ and $A\subseteq \N,$ denote by 

                       \begin{itemize}
                           \item $-y+A=\{n: n+y\in A\};$ and
                           \item $y^{-1}A=\{n: ny\in A\}.$
                       \end{itemize}
    \end{enumerate}
\end{defn}

To prove our result, we need the following three technical results.
The following lemma is a reformulation of the Hindman theorem \cite{finitesum}.
\begin{lem}\textup{\cite[Lemma 5.19.2.]{hs}}\label{1}
    Any finite coloring of an $IP$ set contains a monochromatic IP subset. 
\end{lem}

The proof of the next two lemmas is standard, but we give the proofs here for the sake of completeness.

\begin{lem}\label{2}
For any IP set \( A \subseteq \mathbb{N} \) and any natural number \( y \), there exists another IP set \( A' \subseteq A \) satisfying \( y^{-1}A' \subseteq \mathbb{N} \).
\end{lem}

\begin{proof}
By passing to a subsequence if necessary, we may assume that 
\(A\) contains \(FS(\langle x_n \rangle)\) for some sequence 
\(\langle x_n \rangle \subset (y\mathbb{N}+j)\), where 
\(j \in \{0,1,\ldots ,y-1\}\).

Now let \(F\) be any finite set with \(|F|=y\). Then
\[
x_F=\sum_{n\in F} x_n \in y\mathbb{N}.
\]
Thus every sum of exactly \(y\) distinct terms from the sequence 
is divisible by \(y\).

Selecting infinitely many distinct elements from the sequence 
\(\langle x_F\rangle\) and using them as generators, it follows that 
\(A\) contains another finite sums set
\[
A' = FS(\langle y_n\rangle),
\]
where each \(y_n\) is of the form
\[
y_n = y z_n
\]
for some integer \(z_n\). And hence $$FS(\langle z_n\rangle)= y^{-1}A'\subseteq \mathbb{N}.$$ This completes the proof.
\end{proof}


\begin{lem}\label{3}
For any \( IP \) set \( A \subseteq \mathbb{N} \), there exists \( y \in A \) such that 
\[
(-y + A) \cap A
\]
is also an \( IP \) set.
\end{lem}

\begin{proof}
Let \( A = FS(\langle x_n \rangle_n) \) be an \( IP \) set generated by a sequence \( \langle x_n \rangle_n \), i.e. 
\[
A = \left\{ \sum_{n \in F} x_n :  F \in \mathcal{P}_f(\N) \right\}.
\]
Take \( y = x_1 \in A \). Then any finite sum of distinct elements from \( (x_{n})_{n \ge 2} \) is of the form \( \sum_{n \in F} x_n \) with \( F \subseteq \{2,3,\dots\} \), and hence 
\[
\sum_{n \in F} x_n + y = \sum_{n \in F \cup \{1\}} x_n \in A.
\]
Thus, every element of \( FS(x_{n})_{n \ge 2} \) belongs to \( (-y + A) \cap A \). Since \( FS(x_{n})_{n \ge 2} \) is itself an \( IP \) set, the intersection \( (-y + A) \cap A \) contains an \( IP \) set, and hence is \( IP \).
\end{proof}



Now we are in the position to prove our Theorem \ref{p}.

\begin{proof}[Proof of  Theorem \ref{p}:] Suppose that \( \mathbb{N} \) is finitely colored, say  
\[
\mathbb{N} = \bigcup_{i=1}^r A_i.
\]
Without loss of generality, assume that \( A_1 \) is an \( IP \) set, and set \( D_0 = A_1 \).
By Lemma~\ref{3}, we can choose an element \( y_1 \in A_1 \) such that 
\( (-y_1 + A_1) \cap A_1 \) is an \( IP \) set. 
Then, by Lemma~\ref{2}, we can choose an $IP$ set
\[
C_1 \subset y_1^{-1}\big( (-y_1 + A_1) \cap A_1 \big),
\]
which is a subset of \( \mathbb{N} \). 
Finally, by Lemma~\ref{1}, there exists \( i_1 \in \{1, 2, \ldots, r\} \) such that
\[
D_1 = C_1 \cap A_{i_1}
     = y_1^{-1}\big( (-y_1 + A_1) \cap A_1 \big) \cap A_{i_1}
\]
is an \( IP \) set.





Using Lemmas~\ref{2} and~\ref{3}, choose \( y_2 \in D_1 \) such that 
\[
C_2 \subseteq y_2^{-1}\big( (-y_2 + D_1) \cap D_1 \big) \subset \mathbb{N}
\]
is an \( IP \) set. 
Then, by Lemma~\ref{1}, there exists \( i_2 \in \{1,2,\ldots,r\} \) such that 
\( D_2 = C_2 \cap A_{i_2} \) is an \( IP \) set.
Proceeding inductively, for each \( n \in \mathbb{N} \), choose \( y_n \in D_{n-1} \) and \( i_n \in \{1,2,\ldots,r\} \) such that the set 
    \[
    D_n = y_n^{-1}\big( (-y_n + D_{n-1}) \cap D_{n-1} \big) \cap A_{i_n}
    \]
    is an \( IP \) set.
    



By the Pigeonhole Principle, there exists \( k \in \{1,2,\ldots,r\} \) such that 
\( A_{i_j} = A_{i_n} = A_k \) for some indices \( j < n \).
Now choose \( x \in D_n \subseteq A_{i_n} = A_k \).

From the construction in~(2), we have
\[
x y_n \in D_{n-1} \subseteq y_{n-1}^{-1} D_{n-2} \subseteq \cdots \subseteq 
y_{n-1}^{-1} \cdots y_{j+1}^{-1} D_j.
\]
Hence,
\[
x y_n \cdots y_{j+1} \in D_j \subseteq A_{i_j} = A_k.
\]
Moreover, since \( y_n \in D_{n-1} \), applying the same argument gives
\[
y_n \cdots y_{j+1} \in D_j \subseteq A_{i_j} = A_k.
\]
Thus, both \( y_n \cdots y_{j+1} \) and \( x y_n \cdots y_{j+1} \) lie in the same color class \( A_k \).


Again, observe that 
\[
x y_n + y_n \in D_{n-1},
\] 
so by the same argument as above,
\[
(x y_n + y_n) \, y_{n-1} \cdots y_{j+1} \in D_j \subseteq A_{i_j}.
\]

Now define
\begin{enumerate}
    \item \( b = x \), and
    \item  \( a = y_n \cdots y_{j+1} \).
\end{enumerate}

Then \( a, b \in A_k \), and
\[
x y_n \cdots y_{j+1} = a b \in A_k, \quad \text{and} \quad 
(x y_n + y_n) \, y_{n-1} \cdots y_{j+1} = a(b+1) \in A_k.
\]

Hence, the set
\[
\{a, b, ab, a(b+1)\} \subseteq A_k
\]
is monochromatic. 

This completes the proof.

    
\end{proof}





 %The body of your paper goes here~\cite{cilleruelo}.

%\newpage %% AUTHOR: please comment out this line.  It serves only
%%   to demonstrate both types of header line in aic-template.pdf

%\section{Expansion estimates}

% More of the body of your paper goes here~\cite{bergelson-johnson-moreira}.

%%% AUTHOR: optional appendix here
%\appendix %% you may comment this out if no Appendix
%\section*{Appendix}
%\section{Improving the constants}
%Material is placed here as needed.

%%% AUTHOR: optional acknowledgments here
\section*{Acknowledgments} %%  you may comment this out if no Ackno
The author is thankful to the referees for their comments on the previous draft of this paper and also for the proof of the Lemma \ref{2}. This paper is supported by NBHM postdoctoral fellowship with reference no: 0204/27/(27)/2023/R \& D-II/11927.

%%% AUTHOR:
%%% Bibliography goes here. Note that the arXiv cannot process bibtex
%%% or biber bibliographies.  Example of acceptable bibliograpy format:
\bibliographystyle{amsplain}
\begin{thebibliography}{99}


\bibitem{al1} Ryan Alweiss. \newblock Monochromatic sums and products of polynomials. 
\newblock {\em Discrete Anal}, 2024:5, 7 pp.

\bibitem{erdos} P. Erd\H{o}s: \textit{Problems and results on combinatorial number theory. III}, {Lecture Notes in Math.} {626}, Springer, Berlin, 1977, pp. 43–72. %\href{https://doi.org/10.1007/bfb0063064}{link}.

%\bibitem{.} Ronald L. Graham, Bruce L. Rothschild, and Joel H. Spencer.
%\newblock Ramsey Theory, second ed.
%\newblock {\em Wiley-Interscience Series in Discrete Mathematics and Optimization, John Wiley \& Sons, New York}, 1990.

\bibitem{finitesum} Neil Hindman. 
\newblock Finite sums from sequences within cells of partitions of $\mathbb{N}$.
\newblock {\em J. Combin. Theory Ser. A}, 17 (1974), 1–11.

\bibitem{trans} Neil Hindman. 
\newblock Partitions and sums and products of integers.
\newblock {\em Trans. AMS.}, 247 (1979), 227–245.

\bibitem{hs} Neil Hindman, and Dona Strauss.
\newblock Algebra in the Stone-\v{C}ech Compactifications: theory and applications, second edition.
\newblock {\em de Gruyter}, Berlin, 2012.

\bibitem{m} Joel Moreira.
\newblock Monochromatic sums and products in $\N$.
\newblock {\em Ann. Math.}, (2) 185,  (2017), no. 3, 1069–1090.

\bibitem{saha} Julian Sahasrabudhe.
\newblock Exponential patterns in arithmetic ramsey theory.
\newblock {\em Acta Arith.}, 182 (2017), 13.

\bibitem{schur} Issai Schur.
\newblock \"{U}ber die Kongruenz $x^{m}+y^{m}\equiv z^{m}~(\!\!\!\!\mod p)$.
\newblock {\em Jahresber. Dtsch. Math.-Ver.}, 25 (1916), 114–117.

\end{thebibliography}
%% AUTHOR: You can generate such a bibliography from a .bib file by 
%% running pdflatex/bibtex/pdflatex/pdflatex and then pasting the .bbl file
%% between \begin{thebibliography} and \end{bibliography}

%%% AUTHOR: Include a short description of each author following the
%%% structure below. Use the same short tags used previously.  
%%% Use \imageat{} and \imagedot{} instead of "@" and "." in
%%% email addresses-this replaces the symbols with graphics to avoid 
%%% e-mail address harvesting from the .pdf file
\begin{aicauthors}
\begin{authorinfo}[SG]
  Sayan Goswami\\
  Ramakrishna Mission Vivekananda Educational and Research Institute (RKMVERI), \\
  Belur, India\\
  sayan92m\imageat{}gmail\imagedot{}com \\
  %\url{https://www.cs.elte.hu/erdos}
\end{authorinfo}
\end{aicauthors}

\end{document}
